% ============================================================
%  preamble.tex —— 宏包、全局样式与自定义命令
%  编译：xelatex（需 ctex；建议 xelatex 连续编译两次以生成目录/交叉引用）
% ============================================================

% --- 页面 ---
\usepackage{geometry}
\geometry{a4paper, left=3cm, right=3cm, top=3cm, bottom=3cm}

% --- 数学与符号 ---
\usepackage{amsmath}
\usepackage{amssymb}

% --- 图形与表格 ---
\usepackage{graphicx}
\usepackage{booktabs}
\usepackage{longtable}
\usepackage{tabularx}
\usepackage{array}
\usepackage{multirow}
\usepackage{float}
\usepackage{caption}

% --- 列表 ---
\usepackage{enumitem}
\setlist{nosep, leftmargin=2.2em}

% --- 颜色与代码 ---
\usepackage{xcolor}
\usepackage{listings}

% --- 绘图 ---
\usepackage{tikz}
\usetikzlibrary{positioning, arrows.meta, shapes.geometric, fit, backgrounds, calc, automata}

% --- 超链接（放在较后）---
\usepackage[hidelinks]{hyperref}

% --- 下划线在正文中正常显示（避免大量 \_ 转义）；须在 hyperref 之后加载 ---
\usepackage{underscore}

% --- 行距 ---
\linespread{1.4}

% --- 图表编号按“章-序号”，标签与标题间用空格（中文习惯：图 4-1 标题）---
\renewcommand{\thefigure}{\thechapter-\arabic{figure}}
\renewcommand{\thetable}{\thechapter-\arabic{table}}
\renewcommand{\theequation}{\thechapter-\arabic{equation}}
\captionsetup{labelsep=space, font=small}

% --- 表格列类型 ---
\newcolumntype{Y}{>{\raggedright\arraybackslash}X}
\newcolumntype{C}[1]{>{\centering\arraybackslash}p{#1}}
\newcolumntype{L}[1]{>{\raggedright\arraybackslash}p{#1}}

% --- 代码样式 ---
\definecolor{codebg}{RGB}{248,248,248}
\definecolor{codekw}{RGB}{0,0,170}
\definecolor{codecmt}{RGB}{0,120,0}
\definecolor{codestr}{RGB}{160,0,0}
\definecolor{rulegray}{gray}{0.6}
\lstset{
  basicstyle=\ttfamily\footnotesize,
  backgroundcolor=\color{codebg},
  keywordstyle=\color{codekw}\bfseries,
  commentstyle=\color{codecmt},
  stringstyle=\color{codestr},
  numberstyle=\tiny\color{gray},
  breaklines=true,
  breakatwhitespace=false,
  showstringspaces=false,
  frame=single,
  rulecolor=\color{rulegray},
  columns=fullflexible,
  keepspaces=true,
  tabsize=2,
}

% --- 行内代码 ---
\newcommand{\code}[1]{\texttt{#1}}

% --- 带圈数字（①② 在拉丁/中文字体中可能缺字，统一用此宏绘制）---
\newcommand{\cnum}[1]{\textcircled{\raisebox{0.3pt}{\scriptsize #1}}}

% --- 减少 overfull：允许段落末行轻微伸展 ---
\emergencystretch=2em

% --- TikZ 通用样式 ---
\tikzset{
  box/.style   ={draw, rounded corners, align=center, inner sep=5pt,
                 minimum height=9mm, font=\small, fill=white},
  db/.style    ={draw, cylinder, shape border rotate=90, aspect=0.25,
                 align=center, inner sep=4pt, font=\footnotesize, fill=white},
  zone/.style  ={draw, dashed, rounded corners, inner sep=8mm},
  flow/.style  ={-{Stealth[length=2.5mm]}, thick},
  flowbi/.style={{Stealth[length=2.5mm]}-{Stealth[length=2.5mm]}, thick},
  lbl/.style   ={font=\footnotesize, align=center},
}
